\section{Punctured centering of the regular ultra coefficient}
\label{sec:tpc37-punctured-centering}

After \cref{sec:tpc37-physical-zero-faces}, the physical equality is
restricted to \(q\nmid N_\alpha(j)\).  Hence both factors in
\(su=N_\alpha(j)\) are nonzero modulo \(q\).  We now center the completed
ultra coefficient without reintroducing an artificial value at the zero
residue.

\subsection{An exact asymmetric centering identity}

Let \(q>2\) be prime, put \(G_q=\mathbb F_q^\times\), and fix
\(h\in G_q\).  Let \(F,G:\mathbb F_q\to\mathbb C\), where
\begin{equation}
 G(0)=0.
 \label{eq:tpc37-raw-g-zero}
\end{equation}
The function \(F\) is the raw residue aggregation of a \(d\)-weight, and
\(G\) is the raw aggregation of the regular ultra weight.  Define the
ordinary additive centering of \(F\) by
\begin{equation}
 \overline F=\frac1q\sum_{D\in\mathbb F_q}F(D),
 \qquad f(D)=F(D)-\overline F.
 \label{eq:tpc37-f-centering}
\end{equation}
For \(G\), instead put
\begin{equation}
 \beta=\frac1{q-1}\sum_{U\in G_q}G(U),
 \label{eq:tpc37-punctured-mean}
\end{equation}
and define its punctured centering by
\begin{equation}
 g(0)=0,
 \qquad
 g(U)=G(U)-\beta\quad(U\in G_q).
 \label{eq:tpc37-punctured-g}
\end{equation}
Then
\begin{equation}
 \sum_D f(D)=\sum_U g(U)=0,
 \qquad g(0)=0.
 \label{eq:tpc37-punctured-properties}
\end{equation}

For \(a,s\in G_q\), write
\begin{equation}
 \mathcal C_{a,s}(A,B)
 :=\sum_{D\in\mathbb F_q}
 A(D)B\!\left(s^{-1}(aD+h)\right).
 \label{eq:tpc37-affine-correlation}
\end{equation}

\begin{proposition}[Punctured-centering identity]
\label{prop:tpc37-punctured-identity}
For every \(a,s\in G_q\),
\begin{equation}
 \boxed{
 \mathcal C_{a,s}(F,G)
 =\mathcal C_{a,s}(f,g)
 -\beta f(-h/a)
 +(q-1)\overline F\,\beta.}
 \label{eq:tpc37-punctured-identity}
\end{equation}
\end{proposition}

\begin{proof}
On \(\mathbb F_q\),
\begin{equation}
 F=f+\overline F,
 \qquad
 G=g+\beta\mathbf 1_{G_q}.
 \label{eq:tpc37-punctured-decompositions}
\end{equation}
The affine map \(D\mapsto s^{-1}(aD+h)\) is a permutation, so the
constant part of \(F\) contributes
\begin{equation}
 \overline F\sum_U G(U)=(q-1)\overline F\,\beta.
 \label{eq:tpc37-F-mean-contribution}
\end{equation}
For the other correction, the indicator
\(\mathbf 1_{G_q}(s^{-1}(aD+h))\) deletes exactly
\(D_0=-h/a\).  Since \(f\) is centered,
\begin{equation}
 \sum_D f(D)\mathbf 1_{aD+h\ne0}
 =-f(D_0).
 \label{eq:tpc37-deleted-residue}
\end{equation}
Combining the two identities proves
\eqref{eq:tpc37-punctured-identity}.
\end{proof}

The identity extends directly to the product--slope topology.  Define
\begin{align}
 R_s^{\rm raw}(a)&:=\mathcal C_{a,s}(F,G),
 &R_s(a)&:=\mathcal C_{a,s}(f,g),
 \label{eq:tpc37-raw-centered-R}\\
 (T^{\rm raw}b)(j,s)&:=\sum_{\ell\in G_q}b(\ell)R_s^{\rm raw}(\ell j),
 &(Tb)(j,s)&:=\sum_{\ell\in G_q}b(\ell)R_s(\ell j).
 \label{eq:tpc37-raw-centered-T}
\end{align}
Put
\begin{equation}
 B_0=\sum_{\ell\in G_q}b(\ell),
 \qquad
 H_f(j)=\sum_{\ell\in G_q}
 b(\ell)f\!\left(-\frac h{\ell j}\right).
 \label{eq:tpc37-rank-one-data}
\end{equation}
Then \cref{prop:tpc37-punctured-identity} gives the pointwise formula
\begin{equation}
 \boxed{
 (T^{\rm raw}b)(j,s)
 =(Tb)(j,s)-\beta H_f(j)
 +(q-1)\overline F\,\beta B_0.}
 \label{eq:tpc37-product-punctured-identity}
\end{equation}
For an arbitrary weight \(w:G_q\to\mathbb C\), with
\(W_0=\sum_s w(s)\),
\begin{align}
 \sum_s w(s)(T^{\rm raw}b)(j,s)
 ={}&\sum_s w(s)(Tb)(j,s) \notag\\
 &+W_0\{-\beta H_f(j)
 +(q-1)\overline F\,\beta B_0\}.
 \label{eq:tpc37-weighted-punctured-identity}
\end{align}
Thus every centering correction is explicit and contains the single
scalar \(\beta\).  No smallness assumption on \(\overline F\) is used in
this algebraic statement.

\subsection{The actual M\"obius mean}

On a reciprocal dyadic block, choose the canonical allocation in which
the ultra factor remains
\begin{equation}
 a_U(u)=-\mu(u)\log u\,\Phi_U(u),
 \label{eq:tpc37-clean-ultra-weight}
\end{equation}
where \(\Phi_U\) is the indicator of
\(\{u:U<u\le2U,\ T<u\le U_0\}\), or a bounded smooth/BV cutoff with
controlled seminorms on that interval.  On the regular physical support one may
insert \(\mathbf 1_{q\nmid u}\) exactly.  The raw residue aggregation is
therefore
\begin{equation}
 G_U(r)=
 \sum_{\substack{u\equiv r\pmod q\\q\nmid u}}
 -\mu(u)\log u\,\Phi_U(u),
 \qquad r\in\mathbb F_q,
 \label{eq:tpc37-regular-ultra-periodization}
\end{equation}
and indeed \(G_U(0)=0\).  Its punctured mean is
\begin{equation}
 \beta_U
 =-\frac1{q-1}
 \sum_{\substack{u\ge1\\q\nmid u}}
 \mu(u)\log u\,\Phi_U(u).
 \label{eq:tpc37-physical-beta}
\end{equation}

\begin{proposition}[Uniform one-M\"obius bound for the punctured mean]
\label{prop:tpc37-beta-bound}
Assume \(q\asymp J\), \(U\ge T\), and
\eqref{eq:tpc37-endpoint-scales}, and assume that the cutoff family in
\eqref{eq:tpc37-clean-ultra-weight} has uniformly bounded smooth or BV
seminorms.  For every fixed \(A>0\),
\begin{equation}
 \boxed{
 |\beta_U|
 \ll_{A,\Phi}\frac Uq(\log X)^{-A}.}
 \label{eq:tpc37-beta-bound}
\end{equation}
The estimate is uniform over the dyadic ultra blocks.
\end{proposition}

\begin{proof}
Let
\begin{equation}
 M(x)=\sum_{n\le x}\mu(n),
 \qquad
 M_q(x)=\sum_{\substack{n\le x\\(n,q)=1}}\mu(n).
 \label{eq:tpc37-Mertens-functions}
\end{equation}
Separating the multiples of the prime \(q\) gives the exact recurrence
\begin{equation}
 M_q(x)=\sum_{r\ge0}M(x/q^r),
 \label{eq:tpc37-coprime-Mertens-recurrence}
\end{equation}
where the sum terminates.  The standard zero-free-region form of the
prime number theorem \citep[Chapters~5 and~18]{IwaniecKowalski2004}
gives, for every fixed \(B>0\),
\begin{equation}
 M(y)\ll_B y(\log y)^{-B}
 \label{eq:tpc37-standard-Mertens-bound}
\end{equation}
when \(y\) is large.  In
\eqref{eq:tpc37-coprime-Mertens-recurrence}, the terms after the first
have total trivial size \(O(x/q)\), which is smaller than
\(x(\log X)^{-B}\) because \(q=X^{133/400+o(1)}\).  Hence, uniformly for
\(x\asymp U\),
\begin{equation}
 M_q(x)\ll_B U(\log X)^{-B}.
 \label{eq:tpc37-uniform-coprime-Mertens}
\end{equation}
Partial summation, with \(B\) chosen beyond \(A\) and the fixed seminorm
loss of \(\Phi_U\), yields
\begin{equation}
 \sum_{q\nmid u}\mu(u)\log u\,\Phi_U(u)
 \ll_{A,\Phi}U(\log X)^{-A}.
 \label{eq:tpc37-weighted-one-Mobius}
\end{equation}
Division by \(q-1\asymp q\) proves
\eqref{eq:tpc37-beta-bound}.
\end{proof}

The endpoint separation
\begin{equation}
 \frac Tq=X^{107/2000+o(1)}
 \label{eq:tpc37-ultra-quotient-length}
\end{equation}
also permits the same argument after writing a zero-residue sum as
\(u=qv\).  More importantly, if \(|\mathcal S_U|\ll X/U\), then
\begin{equation}
 |\beta_U|\,|\mathcal S_U|
 \ll_A\frac Xq(\log X)^{-A}
 \asymp Q(\log X)^{-A}.
 \label{eq:tpc37-beta-times-s-length}
\end{equation}
This is an arbitrary logarithmic saving, not a fixed power of \(X\).
Even a Vinogradov--Korobov error has size \(X^{-o(1)}\), so
\cref{prop:tpc37-beta-bound} alone does not pay the
\(X^{1/400}\) endpoint margin.  Nor may its logarithmic saving be placed
after an uncontrolled worst-case \(X^{o(1)}\) absolute reassembly.  The
mean must be extracted while the ultra weight still has the clean form
\eqref{eq:tpc37-clean-ultra-weight}.

\subsection{Exact cancellation of complete residue cycles}

The punctured-centered core has no complete complementary-factor trace.

\begin{proposition}[Complete-cycle annihilation]
\label{prop:tpc37-complete-cycle-annihilation}
For every \(a,j\in G_q\),
\begin{equation}
 \sum_{s\in G_q}R_s(a)=0,
 \qquad
 \sum_{s\in G_q}(Tb)(j,s)=0.
 \label{eq:tpc37-complete-cycle-annihilation}
\end{equation}
\end{proposition}

\begin{proof}
Fix \(a\) and put \(D_0=-h/a\).  The term \(D=D_0\) vanishes because
\(g(0)=0\).  For \(D\ne D_0\), as \(s\) runs over \(G_q\), the value
\(s^{-1}(aD+h)\) runs once over \(G_q\).  Its total contribution is zero
because \(\sum_{U\in G_q}g(U)=0\).  This proves the first identity; the
second follows by summing in \(\ell\).
\end{proof}

For an integer interval \(\mathcal S\), define its nonzero residue lift
count by
\begin{equation}
 W_{\mathcal S}(r)
 :=\#\{s\in\mathcal S:s\equiv r\pmod q\},
 \qquad r\in G_q.
 \label{eq:tpc37-s-lift-count}
\end{equation}
Writing \(|\mathcal S|=kq+r_0\), one has
\begin{equation}
 W_{\mathcal S}(r)=k+\rho_{\mathcal S}(r),
 \qquad
 \rho_{\mathcal S}(r)\in\{0,1\},
 \qquad
 |\operatorname{supp}\rho_{\mathcal S}|<q.
 \label{eq:tpc37-cycle-boundary-decomposition}
\end{equation}
Consequently,
\begin{equation}
 \sum_{r\in G_q}W_{\mathcal S}(r)(Tb)(j,r)
 =\sum_{r\in G_q}\rho_{\mathcal S}(r)(Tb)(j,r).
 \label{eq:tpc37-complete-cycles-removed}
\end{equation}
For the complementary interval \(|\mathcal S_U|\asymp X/U\), the number
of removed complete cycles satisfies
\begin{equation}
 k_U\ll1+\frac X{Uq}\asymp1+\frac Q U.
 \label{eq:tpc37-complete-cycle-count}
\end{equation}
It is \(\asymp Q/U\) whenever the complementary interval contains at
least a fixed positive number of complete residue cycles.  In particular,
\begin{equation}
 k_T=X^{563/2000+o(1)},
 \label{eq:tpc37-low-ultra-cycle-count}
\end{equation}
whereas \(k_U=0\) once \(U\) exceeds a sufficiently large fixed multiple
of \(Q\).

Equation \eqref{eq:tpc37-complete-cycles-removed} concerns the core
residue lift count.  An alias selector or another nonconstant reassembly
weight belongs to \(\rho\)-type weighted data and is not annihilated by
this proposition.

\subsection{Comparison with ordinary additive centering}

If one instead uses
\begin{equation}
 \overline G=\frac1q\sum_U G(U)=\frac{q-1}{q}\beta,
 \qquad
 g^{\rm add}=G-\overline G,
 \label{eq:tpc37-ordinary-g-centering}
\end{equation}
then
\begin{equation}
 g^{\rm add}(0)=-\frac{q-1}{q}\beta
 \label{eq:tpc37-ordinary-centered-zero}
\end{equation}
despite the raw identity \(G(0)=0\).  In fact,
\begin{equation}
 g^{\rm add}
 =g+\frac\beta q-\beta\mathbf 1_{\{0\}}.
 \label{eq:tpc37-two-centerings-comparison}
\end{equation}
The complete-trace formula for ordinary centering becomes
\begin{equation}
 \sum_{s\in G_q}
 \mathcal C_{a,s}(f,g^{\rm add})
 =-(q-1)\beta f(-h/a).
 \label{eq:tpc37-ordinary-complete-trace}
\end{equation}
Thus the ordinary value \(g^{\rm add}(0)\) is a centering boundary, not a
physical coefficient at an allowed zero residue.  The punctured identity
\eqref{eq:tpc37-punctured-identity} records the same quantity as an
explicit correction that is rank one in the complementary
\(s\)-coordinate while keeping the core value \(g(0)=0\).
It is an exact reorganization, not an additional cancellation theorem.

The scope is therefore precise.  The physical zero coordinate faces are
handled by \cref{thm:tpc37-one-prime-face}; punctured centering removes
the artificial zero value and complete constant cycles; and
\cref{prop:tpc37-beta-bound} controls the exposed clean M\"obius mean.
The remaining nonconstant weighted core, the CRT aliases, and the
terminal \(s=1\) layer are not estimated here.  In particular, this
section does not prove the full coefficient gate.
